\documentclass[a4paper,12pt]{article}
\usepackage{graphicx}
\usepackage{caption}
\usepackage{subcaption}
\usepackage{float}
\usepackage[margin=1in]{geometry}
\usepackage{lmodern}
\usepackage{amsmath}
\usepackage{csvsimple}

\begin{document}

\title{Relatório de Adaline e Perceptron}
\author{Eduardo}
\date{\today}
\maketitle

\section*{Adaline}

Os resultados obtidos para cada fold do K-Fold estão apresentados na Tabela~\ref{tab:error_stats}.

\begin{table}[H]
  \centering
  \caption{Resultados do K-Fold}
  \label{tab:error_stats}
  \csvautotabular{./error_stats.csv}
\end{table}

A figura \ref{fig:fold_predictions} compara as previsões do Adaline com os valores reais para cada fold do K-Fold. 

\begin{figure}[H]
  \centering
  \includegraphics[width=0.8\textwidth]{./fold_predictions.png}
  \caption{Comparação entre previsões e valores reais do Adaline}
  \label{fig:fold_predictions}
\end{figure}

\section*{Conclusão Adaline}

Por sua simplicidade, o modelo apresenta resultados satisfatórios para este conjunto de dados. Entretanto, redes com mais camadas podem ser mais adequadas para capturar padrões complexos.

\section*{Perceptron}

Os resultados obtidos para cada fold do K-Fold estão apresentados na Tabela~\ref{tab:kfold_results}.

\begin{table}[H]
  \centering
  \caption{Resultados do K-Fold}
  \label{tab:kfold_results}
  \csvautotabular{./kfold_results.csv}
\end{table}

\section*{Conclusão Perceptron}

Os resultados obtidos demonstram que o modelo Perceptron apresentou uma precisão média de 83\% ao longo das 10 iterações do K-Fold, com um desvio padrão de 14\%. Isso indica que o modelo possui um desempenho consistente, mas com alguma variação entre os folds. A alta precisão em alguns folds, como 95\%, sugere que o modelo é capaz de capturar bem os padrões em determinados subconjuntos dos dados. No entanto, a menor precisão de 42\% em um dos folds aponta para possíveis limitações do modelo em lidar com certos cenários ou distribuições dos dados. Ajustes nos hiperparâmetros ou técnicas adicionais de pré-processamento podem ser explorados para melhorar a robustez do modelo.

\end{document}